\chapter{Textual IR}
\label{ch:pebble-text}

When a compiler fails, the smallest useful artifact is often the IR just before the failing pass.
A printer that emits attractive but unparsable annotations is a diagnostic aid, not a reproducible interchange format.
PebbleIR's first format favors a stable contract over source-language elegance.

\section{A versioned record format}
\label{sec:pebble-records}
The header contains the magic word, format version, block count, and allocated value count.
Value records follow in ascending ID order, including tombstones.
Each record contains ID, owning block, opcode name, type name, operands $a,b$, immediate, and predecessor IDs $p,q$.
Terminator records follow in block order, then a final \code{end} line.
Every record ends in a newline, including the final record; trailing records are rejected.

Listing~\ref{lst:pebble-text-example} is the unoptimized running expression.
All zeros are meaningful canonical unused fields.
The parser rejects a nonzero predecessor field on an addition, even though the interpreter would not read that field.

\begin{irlisting}[caption={The version-one format preserves value IDs and explicitly encodes every field.},label={lst:pebble-text-example}]
pebble 1 1 7
v 1 0 input u8 0 0 0 0 0
v 2 0 const u8 0 0 7 0 0
v 3 0 add u8 1 2 0 0 0
v 4 0 const u8 0 0 3 0 0
v 5 0 const u8 0 0 5 0 0
v 6 0 add u8 4 5 0 0 0
v 7 0 mul u8 3 6 0 0 0
t 0 ret 7 0 0
end
\end{irlisting}

\section{Parsing is a transaction}
\label{sec:pebble-parse-transaction}
The reader accepts only bounded line records, recognized names, and unsigned decimal fields with checked accumulation.
It does not use unchecked \code{atoi}, whose interface does not provide the validation required here.
It constructs a temporary function, checks the complete syntax, calls the verifier, and only then assigns the caller's destination.
On failure the destination remains unchanged, although the input stream has advanced.
Algorithm~\ref{alg:pebble-text-parse} states that separation.

\begin{algorithm}[htbp]
\caption{Parsing a complete IR file without publishing a partial function.}
\label{alg:pebble-text-parse}
\begin{pseudocode}
\Require readable stream $S$, destination $F$
\Ensure verified new lifetime in $F$, or unchanged $F$
\State parse version and bounded counts into temporary $T$
\State read exactly the declared value records in ID order
\State read exactly the declared terminators in block order
\State require final marker, newline, and end of stream without error
\State verify $T$; reject if any structural or semantic invariant fails
\State publish $T$ and discard facts from the old lifetime of $F$
\end{pseudocode}
\end{algorithm}

For an input of $C$ characters, lexical processing is $O(C)$ before verification.
The 256-character line buffer is a format implementation limit: an overlong record is rejected, not split into several records.
The source-expression parser is a different client and uses different grammar and size limits.
Do not accidentally accept a source expression where an IR snapshot is expected merely because both are text.

\section{What round trips promise}
\label{sec:pebble-roundtrip}
The test prints a verified function, parses it, prints the parsed function, and compares the two byte streams.
It performs this check both before and after optimization, including loops and tombstones.
The property is canonical print stability, not byte-for-byte preservation of arbitrary input whitespace.
Parsing may accept additional whitespace between tokens and printing removes that variation.

Revisions and C addresses are not serialized.
They describe a live editor object, not program meaning.
A reparsed function gets a new logical lifetime with the same serialized IDs and semantics.
Debug source spans, pass names, and failure inputs could later be stored in a separate reproducer envelope so they do not change the core language's identity rules.

\section{Exercises}
\label{sec:pebble-seven-exercises}
\begin{enumerate}
\item \label{ex:pebble-text-fold} Rewrite Listing~\ref{lst:pebble-text-example} after folding and dead-code elimination. Preserve all seven allocated IDs.
\item \label{ex:pebble-roundtrip-contract} Does print--parse--print stability prove that the interpreter is correct? Give a fault that could pass this test.
\end{enumerate}

\section{Further reading}
\label{sec:pebble-seven-reading}
The C11 stream and conversion interfaces are worth reading alongside a text reader's error contract, especially the distinction between end of file and stream error~\cite[Section~7.21]{wg14-2011-c11}.
